\documentclass[11pt]{article}
\usepackage[margin=1in]{geometry}
\usepackage{amsmath,amssymb,booktabs,hyperref}
\usepackage[numbers]{natbib}
\newcommand{\todo}[1]{\textcolor{red}{[TODO: #1]}}
\usepackage{xcolor}
% RULE: every number in this paper is \input from generated/ (scripts/make_tables.py).
\title{Interpreting World Models: A Tutorial and Methodology}
\author{\todo{authors}}
\begin{document}
\maketitle

\begin{abstract}
A practical guide to making \emph{checkable} claims about what learned world models represent and how they use it.
Every method is paired with the control that would falsify it, and every example runs from the accompanying
\texttt{wmi} toolkit on a testbed with exact ground truth.
\end{abstract}

\section{What counts as a claim about a world model}
Decodability (information is present), use (the model's computation depends on it), and fidelity
(the represented dynamics match the environment). These require different evidence.

\section{Testbeds with ground truth}
KeyGrid: position and key status are exactly known. The POMDP stream shows a full view only at $t=0$, so
position and key status at $t>0$ must be carried in memory---decoding them is non-trivial, unlike
fully observed one-hot inputs where everything is linearly decodable from the raw observation.

\section{Probing without fooling yourself}
\begin{enumerate}
\item Split by episode; assert disjointness (\texttt{wmi.harness.check\_split}).
\item Select probe hyper-parameters on validation only.
\item Always report: raw-input probe, untrained-model probe, shuffled labels, majority class, and a
      control task \citep{hewitt2019control}; report selectivity.
\item Cluster-bootstrap CIs over episodes, not timesteps.
\end{enumerate}
\begin{table}[h]\centering\scriptsize
\IfFileExists{generated/probe_table.tex}{% generated by scripts/make_tables.py; seeds=[0, 1, 2]; commits=['af788c4']; dirty=False
\begin{tabular}{lllccccc}
\toprule
Model & Site & Target & Real & Control task & Selectivity & Shuffled & Majority \\
\midrule
trained & enc & has\_key & 0.649 $\pm$ 0.030 & 0.715 $\pm$ 0.033 & -0.065 $\pm$ 0.018 & 0.611 $\pm$ 0.054 & 0.610 $\pm$ 0.055 \\
trained & enc & x & 0.683 $\pm$ 0.012 & 0.442 $\pm$ 0.035 & 0.241 $\pm$ 0.047 & 0.183 $\pm$ 0.039 & 0.155 $\pm$ 0.012 \\
trained & enc & y & 0.597 $\pm$ 0.008 & 0.440 $\pm$ 0.024 & 0.157 $\pm$ 0.019 & 0.190 $\pm$ 0.047 & 0.164 $\pm$ 0.004 \\
trained & h & has\_key & 0.987 $\pm$ 0.011 & 0.769 $\pm$ 0.055 & 0.218 $\pm$ 0.058 & 0.612 $\pm$ 0.054 & 0.610 $\pm$ 0.055 \\
trained & h & x & 1.000 $\pm$ 0.000 & 0.888 $\pm$ 0.047 & 0.112 $\pm$ 0.047 & 0.145 $\pm$ 0.030 & 0.155 $\pm$ 0.012 \\
trained & h & y & 1.000 $\pm$ 0.000 & 0.841 $\pm$ 0.024 & 0.159 $\pm$ 0.024 & 0.202 $\pm$ 0.017 & 0.164 $\pm$ 0.004 \\
trained & obs & has\_key & 0.646 $\pm$ 0.024 & 0.675 $\pm$ 0.062 & -0.029 $\pm$ 0.085 & 0.613 $\pm$ 0.055 & 0.610 $\pm$ 0.055 \\
trained & obs & x & 0.683 $\pm$ 0.012 & 0.414 $\pm$ 0.021 & 0.269 $\pm$ 0.021 & 0.080 $\pm$ 0.077 & 0.155 $\pm$ 0.012 \\
trained & obs & y & 0.597 $\pm$ 0.008 & 0.387 $\pm$ 0.024 & 0.210 $\pm$ 0.023 & 0.170 $\pm$ 0.048 & 0.164 $\pm$ 0.004 \\
untrained & enc & has\_key & 0.649 $\pm$ 0.030 & 0.685 $\pm$ 0.059 & -0.036 $\pm$ 0.086 & 0.612 $\pm$ 0.053 & 0.610 $\pm$ 0.055 \\
untrained & enc & x & 0.683 $\pm$ 0.012 & 0.471 $\pm$ 0.034 & 0.212 $\pm$ 0.023 & 0.166 $\pm$ 0.112 & 0.155 $\pm$ 0.012 \\
untrained & enc & y & 0.597 $\pm$ 0.008 & 0.410 $\pm$ 0.025 & 0.187 $\pm$ 0.023 & 0.197 $\pm$ 0.076 & 0.164 $\pm$ 0.004 \\
untrained & h & has\_key & 0.701 $\pm$ 0.006 & 0.700 $\pm$ 0.063 & 0.001 $\pm$ 0.062 & 0.612 $\pm$ 0.053 & 0.610 $\pm$ 0.055 \\
untrained & h & x & 0.811 $\pm$ 0.006 & 0.454 $\pm$ 0.038 & 0.358 $\pm$ 0.035 & 0.153 $\pm$ 0.032 & 0.155 $\pm$ 0.012 \\
untrained & h & y & 0.812 $\pm$ 0.007 & 0.445 $\pm$ 0.015 & 0.367 $\pm$ 0.018 & 0.153 $\pm$ 0.034 & 0.164 $\pm$ 0.004 \\
\bottomrule
\end{tabular}
}{\todo{run scripts/run\_probes.py}}
\caption{Probe accuracy on held-out episodes ($t\ge1$). Mean $\pm$ s.d. over seeds.}
\end{table}

\section{From decodability to use: interventions}
Minimal-norm edits that set a probe's logit to a target, compared to matched-norm random edits; measure the
intended effect and a specificity metric on an unrelated output channel.
\begin{table}[h]\centering\scriptsize
\IfFileExists{generated/intervention_table.tex}{
% generated by scripts/make_tables.py; seeds=[0, 1, 2]; commits=['af788c4']; dirty=False
\begin{tabular}{lllcc}
\toprule
Intervention & Direction & $\alpha$ & $p(\text{key at cell})$ & Blocked-channel acc. \\
\midrule
insert\_key\_belief & probe & 0 & 0.000 $\pm$ 0.000 & 1.000 $\pm$ 0.000 \\
insert\_key\_belief & probe & 0.25 & 0.000 $\pm$ 0.000 & 1.000 $\pm$ 0.000 \\
insert\_key\_belief & probe & 0.5 & 0.000 $\pm$ 0.000 & 1.000 $\pm$ 0.000 \\
insert\_key\_belief & probe & 1 & 0.000 $\pm$ 0.000 & 1.000 $\pm$ 0.000 \\
insert\_key\_belief & probe & 2 & 0.001 $\pm$ 0.001 & 1.000 $\pm$ 0.000 \\
insert\_key\_belief & probe & 4 & 0.001 $\pm$ 0.002 & 1.000 $\pm$ 0.000 \\
insert\_key\_belief & diff\_means & 0 & 0.000 $\pm$ 0.000 & 1.000 $\pm$ 0.000 \\
insert\_key\_belief & diff\_means & 0.25 & 0.001 $\pm$ 0.001 & 1.000 $\pm$ 0.000 \\
insert\_key\_belief & diff\_means & 0.5 & 0.004 $\pm$ 0.003 & 1.000 $\pm$ 0.000 \\
insert\_key\_belief & diff\_means & 1 & 0.025 $\pm$ 0.028 & 1.000 $\pm$ 0.000 \\
insert\_key\_belief & diff\_means & 2 & 0.058 $\pm$ 0.038 & 0.986 $\pm$ 0.009 \\
insert\_key\_belief & diff\_means & 4 & 0.105 $\pm$ 0.071 & 0.931 $\pm$ 0.022 \\
insert\_key\_belief & random & 0 & 0.000 $\pm$ 0.000 & 1.000 $\pm$ 0.000 \\
insert\_key\_belief & random & 0.25 & 0.000 $\pm$ 0.000 & 1.000 $\pm$ 0.000 \\
insert\_key\_belief & random & 0.5 & 0.000 $\pm$ 0.000 & 1.000 $\pm$ 0.000 \\
insert\_key\_belief & random & 1 & 0.001 $\pm$ 0.000 & 1.000 $\pm$ 0.000 \\
insert\_key\_belief & random & 2 & 0.004 $\pm$ 0.001 & 0.987 $\pm$ 0.005 \\
insert\_key\_belief & random & 4 & 0.026 $\pm$ 0.006 & 0.931 $\pm$ 0.014 \\
\midrule
remove\_key\_belief & probe & 0 & 0.992 $\pm$ 0.005 & 1.000 $\pm$ 0.000 \\
remove\_key\_belief & probe & 0.25 & 0.992 $\pm$ 0.006 & 1.000 $\pm$ 0.000 \\
remove\_key\_belief & probe & 0.5 & 0.993 $\pm$ 0.007 & 1.000 $\pm$ 0.000 \\
remove\_key\_belief & probe & 1 & 0.992 $\pm$ 0.008 & 1.000 $\pm$ 0.000 \\
remove\_key\_belief & probe & 2 & 0.990 $\pm$ 0.014 & 1.000 $\pm$ 0.000 \\
remove\_key\_belief & probe & 4 & 0.967 $\pm$ 0.054 & 1.000 $\pm$ 0.000 \\
remove\_key\_belief & diff\_means & 0 & 0.992 $\pm$ 0.005 & 1.000 $\pm$ 0.000 \\
remove\_key\_belief & diff\_means & 0.25 & 0.984 $\pm$ 0.008 & 1.000 $\pm$ 0.000 \\
remove\_key\_belief & diff\_means & 0.5 & 0.957 $\pm$ 0.025 & 1.000 $\pm$ 0.000 \\
remove\_key\_belief & diff\_means & 1 & 0.762 $\pm$ 0.122 & 1.000 $\pm$ 0.000 \\
remove\_key\_belief & diff\_means & 2 & 0.429 $\pm$ 0.137 & 0.992 $\pm$ 0.004 \\
remove\_key\_belief & diff\_means & 4 & 0.147 $\pm$ 0.102 & 0.952 $\pm$ 0.014 \\
remove\_key\_belief & random & 0 & 0.992 $\pm$ 0.005 & 1.000 $\pm$ 0.000 \\
remove\_key\_belief & random & 0.25 & 0.992 $\pm$ 0.004 & 1.000 $\pm$ 0.000 \\
remove\_key\_belief & random & 0.5 & 0.990 $\pm$ 0.005 & 1.000 $\pm$ 0.000 \\
remove\_key\_belief & random & 1 & 0.973 $\pm$ 0.004 & 1.000 $\pm$ 0.001 \\
remove\_key\_belief & random & 2 & 0.871 $\pm$ 0.036 & 0.991 $\pm$ 0.006 \\
remove\_key\_belief & random & 4 & 0.591 $\pm$ 0.057 & 0.939 $\pm$ 0.012 \\
\midrule
\multicolumn{5}{l}{$D$ = 1.428 $\pm$ 0.172; cos(probe, diff-means) = 0.033 $\pm$ 0.025} \\
\bottomrule
\end{tabular}
}{\todo{run scripts/run\_probes.py}}
\caption{Editing the has-key direction of the recurrent state.}
\end{table}

\section{Interpretability for planning}
\todo{Latent geometry vs.\ reachability; link to the latent hindsight planning paper.}

\section{Checklist}
\todo{One-page checklist distilled from Sections 3--5.}

\section{Related work}
\citep{li2023othello,probing_rl_wm_2603_21546,wmi_survey_2610_25296,onelens_2606_09936,steer_wam_2607_14943,world2vec_2610_40637}
\todo{Position: this tutorial's contribution is executable methodology with falsification controls.}

\bibliographystyle{plainnat}
\bibliography{../refs}
\end{document}
